% smoke-test.tex — exercises every .cls and commonheader.sty feature.
% Compile with: latexmk -pdf smoke-test.tex
% Verifies Step 2 of the tutor-handouts plan.
\documentclass[lecture]{handout}

\coursename{CS-TEST Skill Smoke Test}
\courseterm{Test Term 2026}
\handoutnumber{1}
\handouttitle{LaTeX Class Smoke Test}

\begin{document}
\maketitle

\section{Typography}
This sentence exercises body text with \code{cudaMemcpy} as inline code and a code listing below.
The Hopper SM90 architecture provides 132 SMs at 1.83 GHz.

\begin{lstlisting}[language=CUDA, caption={Vector add kernel — see CUDA C++ Programming Guide \S 2.2}]
__global__ void vector_add(const float* A, const float* B, float* C, int N) {
    int idx = blockIdx.x * blockDim.x + threadIdx.x;
    if (idx < N) C[idx] = A[idx] + B[idx];
}
\end{lstlisting}

\section{Math}
Hopper's TMA can transfer up to \(2^{16}\) bytes per descriptor:
\[
	\text{throughput} = \frac{\text{tile\_bytes}}{\text{latency} + \text{transfer\_time}}.
\]

\section{Lists}
\begin{itemize}
	\item One concrete number per page (1.83 GHz)
	\item One primary citation (CUDA C++ Programming Guide \S 2.2)
	\item One code listing
\end{itemize}

\section{Tables}
\begin{center}
	\begin{tabular}{lrr}
		\toprule
		\textbf{GPU} & \textbf{SMs} & \textbf{HBM (GB)} \\
		\midrule
		A100         & 108          & 80                \\
		H100         & 132          & 80                \\
		B200         & 148          & 192               \\
		\bottomrule
	\end{tabular}
\end{center}

\section{Triton Listing}
\begin{lstlisting}[language=Triton, caption={Triton vector add — tutorial 01-vector-add.py}]
@triton.jit
def vector_add_kernel(A_ptr, B_ptr, C_ptr, N, BLOCK_SIZE: tl.constexpr):
    pid = tl.program_id(0)
    offsets = pid * BLOCK_SIZE + tl.arange(0, BLOCK_SIZE)
    mask = offsets < N
    a = tl.load(A_ptr + offsets, mask=mask)
    b = tl.load(B_ptr + offsets, mask=mask)
    tl.store(C_ptr + offsets, a + b, mask=mask)
\end{lstlisting}

\end{document}
